\section*{ОПРЕДЕЛЕНИЯ, ОБОЗНАЧЕНИЯ И СОКРАЩЕНИЯ}
\addcontentsline{toc}{section}{ОПРЕДЕЛЕНИЯ, ОБОЗНАЧЕНИЯ И СОКРАЩЕНИЯ}

\begin{enumerate}
	\item Кредитоспособность -- способность заемщика своевременно и в полном объеме выполнять свои кредитные обязательства.

	\item Кредитный риск -- вероятность невыполнения заемщиком обязательств по кредиту.

	\item Дефолт (от англ. Default) -- ситуация, при которой заемщик не выполняет свои обязательства по возврату кредита.

	\item Машинное обучение (от англ. Machine Learning, ML) -- раздел искусственного интеллекта, изучающий методы построения моделей на основе данных.

	\item Градиентный бустинг деревьев (от англ. Gradient Boosting Decision Trees, GBDT) -- ансамблевый метод машинного обучения, основанный на последовательном построении деревьев решений.

	\item Логистическая регрессия (от англ. Logistic Regression) -- статистический метод, используемый для решения задач бинарной классификации.

	\item Метод опорных векторов (от англ. Support Vector Machine, SVM) -- алгоритм машинного обучения для задач классификации и регрессии.

	\item Признак -- характеристика объекта, используемая в модели машинного обучения.

	\item Целевая переменная (от англ. Target Variable) -- переменная, значение которой требуется предсказать.

	\item Обучающая выборка (от англ. Training Dataset) -- набор данных, используемый для обучения модели.

	\item Переобучение (от англ. Overfitting) -- ситуация, при которой модель хорошо работает на обучающих данных, но плохо обобщает на новых данных.

	\item Скоринговый балл -- числовая оценка уровня кредитного риска заемщика.
\end{enumerate}